\section{Scanner Generator-Based dengan Flex}
\label{sec:ch07-flex}

\subsection{Struktur Berkas Flex}

Berkas Flex (\file{.l}) terdiri atas tiga bagian yang dipisahkan
\code{\%\%}:
\begin{enumerate}[label=\arabic*.,leftmargin=*]
  \item \textbf{definisi}: deklarasi, makro, dan kode C;
  \item \textbf{aturan}: pasangan pattern dan aksi;
  \item \textbf{kode pengguna}: fungsi tambahan dan \code{main}.
\end{enumerate}

Flex mengubah aturan menjadi DFA dan menghasilkan fungsi
\code{yylex()} dalam berkas C. Pattern ditulis dengan ekspresi
reguler, sedangkan aksi ditulis dalam C.

\subsection{Definisi dan Makro}

Bagian definisi mendefinisikan makro seperti \code{DIGIT} dan
\code{ID}. Makro ini membuat aturan lebih ringkas dan konsisten.
Potongan bagian definisi dan aturan kata kunci ditunjukkan pada
\lstref{lst:ch07-flex-def}.

\includecode[
  language=Flex,
  firstline=28,
  lastline=57,
  firstnumber=28,
  caption={Definisi dan aturan kata kunci pada \file{lexer\_flex.l}},
  label={lst:ch07-flex-def}
]{../../code/subsetc/src/lexer_flex.l}

Aturan operator dua karakter ditempatkan sebelum operator satu
karakter. Flex memilih pattern terpanjang; jika panjangnya sama,
aturan yang lebih awal menang \citep{flexmanual}. Urutan ini penting untuk memastikan
\code{<=} tidak dipecah menjadi \code{<} dan \code{=}.

\subsection{Aturan Literal dan Identifier}

Aturan untuk literal, identifier, dan karakter ilegal ditunjukkan
pada \lstref{lst:ch07-flex-rules}.

\includecode[
  language=Flex,
  firstline=58,
  lastline=85,
  firstnumber=58,
  caption={Aturan operator, literal, identifier, dan kesalahan},
  label={lst:ch07-flex-rules}
]{../../code/subsetc/src/lexer_flex.l}

Aksi dapat berupa \code{return TOKEN} atau blok C. Pada lexer ini,
aksi hanya mengembalikan jenis token. Nilai literal dan lexeme
diambil dari \code{yytext} oleh driver.

\subsection{Membangun Lexer Flex}

Langkah build:
\begin{lstlisting}[language=bash]
flex -o build/lexer_flex.c src/lexer_flex.l
gcc -Iinclude build/lexer_flex.c src/token.c src/error.c \
    -o build/lexer_flex
\end{lstlisting}

Perintah pertama menghasilkan kode C dari spesifikasi Flex. Perintah
kedua mengompilasi kode tersebut bersama modul token dan error.
Langkah ini memerlukan Flex dan GCC yang terpasang.

\subsection{Perbandingan dengan Lexer Tangan}

\begin{table}[H]
\centering
\caption{Perbandingan lexer tangan dan lexer Flex}
\label{tab:ch07-perbandingan}
\begin{tabularx}{\textwidth}{@{}lXX@{}}
\toprule
\textbf{Aspek} & \textbf{Lexer tangan} & \textbf{Lexer Flex} \\
\midrule
Kontrol & Penuh atas setiap detail & Terbatas pada aksi \\
Keterbacaan & Panjang, tetapi eksplisit & Ringkas dan deklaratif \\
Kualitas kesalahan & Dapat disesuaikan & Perlu kode tambahan \\
Perubahan pattern & Mengubah kode dan menguji ulang & Mengubah aturan \\
Ketergantungan & C standar & Flex dan generator \\
Kinerja & Bergantung implementasi & DFA yang dioptimalkan \\
\bottomrule
\end{tabularx}
\end{table}

Lexer tangan cocok untuk pembelajaran dan sistem yang membutuhkan
diagnostik sangat spesifik. Lexer Flex cocok untuk prototipe cepat
dan bahasa dengan banyak token. Pada proyek Subset-C, keduanya
diimplementasikan agar mahasiswa memahami trade-off-nya.
